
%\needspace{4\baselineskip}
\color{uniblue}{\subsection{Управление выключателем (УВ)}}
\color{black}

\begin{enumerate}[label=\arabic{section}.\arabic{subsection}.\arabic{enumi}, labelsep=4pt, leftmargin=0pt, itemindent=57pt, itemsep=0pt, parsep=0pt] % если что можно крутить itemsep

\item
Функциональный блок управления выключателем предназначен для подачи сигнала на включение силового выключателя. 

\item
В случае присоединения линии к системе через два выключателя для каждого из них можно задействовать свою логику управления выключателем. Далее приведено описание логики управления выключателем В1, логика управления вторым выключателем аналогична.

\item
Функциональная схема логики управления В1 приведена на рисунке \ref{pict_uv:switch_control}. Функциональная схема логики управления В2 выполняется аналогично схеме логики управления В1.

\vspace{2mm}
\begin{figure}[H]
\centering
\includegraphics[width=0.9\textwidth,height=0.9\textheight,keepaspectratio]{\fbpath/761.3391 УВ/img/switch_control.pdf}
\captionof{figure}{Функциональная схема управления выключателем В1}\label{pict_uv:switch_control}
\end{figure}

\item
Команда на включение выключателя «Включить В1» формируется:
\begin{itemize}
    \item от команды оперативного включения выключателя (сигнал \textit{«РКВ В1»});
    \item от схемы полуавтоматического включения (сигнал \textit{«Ср. ПАВ В1»});
    \item при срабатывании АПВ шин или линии (сигналы \textit{«Ср. АПВш В1»} и \textit{«Ср. АПВл В1»}).
\end{itemize}

\item
Оперативное включение выключателя осуществляется с контролем наличия разрешающего сигнала от функционального блока КСН (принцип работы которого описан в пункте 14 настоящего руководства по эксплуатации). Максимальное время, в течение которого продолжается ожидание условий оперативного включения, регулируется уставкой \textbf{«Tожид усл. опер. вкл. В1»}. По истечению данной выдержки времени команда оперативного включения сбрасывается. Выдержка времени на возврат \textbf{«Тпрод опер. вкл. В1»} задает время продления сигнала оперативного включения.

\item
При необходимости может быть введен подхват команды на включение от токового реле, установленного в цепи ЭМВ, в этом случае сигнал на включение удерживается в сработанном состоянии до исчезновения тока через ЭМВ (пока блок-контакты выключателя не разорвут цепь включения). Ввод осуществляется программным ключом \textbf{«Подхват вкл. В1 от РТ ЭМВ»}. 

\item
Предусмотрена возможность блокировки многократных включений на КЗ. Если при наличии сигнала включения будет протекать ток по электромагнитам отключения (что соответствует включению на КЗ), то сигнал \textit{«Включить В1»} заблокируется, пока не пропадут сигналы включения (от АПВ или оперативного включения). Блокировка вводится программным ключом \textbf{«Блок. от многокр. вкл. на КЗ В1»}.

\item
При наличии сигнала \textit{«Блок. включения В1»}, формируемого функциональным блоком контроля силового выключателя, включение выключателя блокируется.

\item
Уставки схемы управления выключателем В1 представлены в таблице \ref{desc_uv:tbl1}. Уставки схемы управления выключателем В2 аналогичны уставкам схемы управления выключателем В1.

\small
\begin{longtable}{|>{\centering\arraybackslash}m{5.3cm}|>{\centering\arraybackslash}m{3.3cm}|>{\centering\arraybackslash}m{4.2cm}|>{\centering\arraybackslash}m{1.8cm}|>{\centering\arraybackslash}m{1cm}|}
\caption{Параметры для настройки схемы управления выключателем В1\hfill\vspace{-0.5\baselineskip}}\label{desc_uv:tbl1}\\ 
\hline
\rowcolor{gray!30}
Параметр & Условное обозначение на схеме & Значение/ Диапазон & Единица измерения & Шаг \\ 
\hline
\endfirsthead
\caption*{\hspace{3pt}\emph{Продолжение таблицы \ref{desc_uv:tbl1}\hfill\vspace{-0.5\baselineskip}}} \\ % сделано по ГОСТ 2.105 п.6.8.7
\hline
\rowcolor{gray!30}
Параметр & Условное обозначение на схеме & Значение/ Диапазон & Единица измерения & Шаг \\ 
%\hline
\endhead

%===m>УВ В1|УВ
\raggedright \textcolor{red}{FullDescription НЕ НАЙДЕН!} & \centering УВ & \centering Введено \\ Выведено & \centering --- & \centering \arraybackslash --- \\
\hline
\raggedright \textcolor{red}{FullDescription НЕ НАЙДЕН!} & \centering Блок. от многокр. вкл. на КЗ & \centering Введено \\ Выведено & \centering --- & \centering \arraybackslash --- \\
\hline
\raggedright \textcolor{red}{FullDescription НЕ НАЙДЕН!} & \centering Подхват вкл. от РТ ЭМВ & \centering Введено \\ Выведено & \centering --- & \centering \arraybackslash --- \\
\hline
\raggedright \textcolor{red}{FullDescription НЕ НАЙДЕН!} & \centering Тожид усл. опер. вкл. & \centering 0 ... 300000 & \centering мс & \centering \arraybackslash 5 \\
\hline
\raggedright \textcolor{red}{FullDescription НЕ НАЙДЕН!} & \centering Тпрод опер. вкл. & \centering 0 ... 300000 & \centering мс & \centering \arraybackslash 5 \\
\hline
%===m
\end{longtable}
\normalsize

\end{enumerate}